\documentclass[11pt,a4paper]{article}

\usepackage[utf8]{inputenc}
\usepackage{amsmath, amssymb, amsthm}
\usepackage{geometry}
\usepackage{hyperref}
\usepackage{mathrsfs}
\usepackage{graphicx}
\usepackage{xcolor}

\geometry{margin=1in}

\title{\Large \textbf{The Circular Convolution of Prime Constellations\\ in Congruence Phase Space}}
\author{Ruqing Chen\\
\small GUT Geoservice Inc., Montreal, Quebec, Canada\\
\small \texttt{ruqing@hotmail.com}}
\date{\today}

\newtheorem{theorem}{Theorem}
\newtheorem{definition}{Definition}
\newtheorem{lemma}{Lemma}

\begin{document}

\maketitle

\begin{abstract}
Following the empirical establishment of the $m$-dimensional proper-depth decay law in Volume I, this paper initiates the theoretical framework of high-dimensional congruence phase-space dynamics. We transition from discrete sieve algorithms to continuous geometric measure theory. By defining the `dead(q)' congruence mechanism as a Castellated Indicator Function on the continuous phase-space torus $\mathbb{T}^1$, we demonstrate that the spatial entanglement of adjacent prime constellations is governed by continuous circular convolution integrals. For Twin-Twin entanglement, this integral evaluates to exactly three quantized interference states. The limit case $j=1$ (adjacent twins merging into a quadruplet) is pure destructive interference, evaluating to $R(1)=\mathfrak{S}_{\text{quad}}/\mathfrak{S}_{\text{twin}}^2\approx0.40$: this quantifies the suppression of the quadruplet \emph{prefactor} (its singular series relative to two independent twins) and is complementary to---not a restatement of---the $-m$ proper-depth slope of Volume~I, which is fixed by the member count. The two effects live on orthogonal axes: interference sets the constant, dimension sets the decay rate.
\end{abstract}

\section{Introduction}
In the first nineteen parts of this series, the Arithmetic Geodynamics framework relied heavily on high-resolution empirical sieving across decimal shells (up to $S_{11}$). The core finding was the rigid scaling law: $d\ln\rho_m / d\ell \to -m$, where $\ell = \ln\ln X$ is the proper depth and $m$ is the tuple dimension~\cite{ChenXVII,ChenXIX,Review}. 

However, to compute the spatial correlation of the 6N lattice---specifically the entanglement between a constellation at center $N$ and another at $N+j$---discrete sieving becomes computationally intractable and theoretically opaque. In this paper, we map the discrete modular arithmetic of the 6N skeleton onto a continuous geometric manifold. We show that prime constellation survival can be modeled as the integration of geometric filters over a phase-space torus $\mathbb{T}^1$.

We stress at the outset that this \emph{horizontal} axis (the lag $j$ between centers) is orthogonal to the \emph{vertical} axis of Volume~I (the proper depth $\ell$). The depth flow $d\ln\rho_m/d\ell\to-m$ is a statement about the $(\ln X)^{-m}$ scaling, fixed by the number of members; the horizontal correlation studied here is a statement about the $X$-independent singular-series constant. The former is a slope, the latter an intercept, and the present paper computes the latter.

\section{The Castellated Indicator Function on $\mathbb{T}^1$}

In the classical circle method~\cite{HL}, the distribution of integers is analyzed on the unit circle. Here, we map the local congruence condition modulo $q$ ($q>3$, prime) to a continuous phase variable $\theta \in [0, 1)$.

\begin{definition}[Castellated Indicator Function]
For a given prime $q$ and a constellation $A$, the forbidden residues (the `dead(q)' mechanism) dictate topological exclusions. We define the Castellated Indicator Function $\Psi_{A,q}(\theta)$ as:
\begin{equation}
\Psi_{A,q}(\theta) = 
\begin{cases} 
0, & \text{if } \theta \in \bigcup_{r \in \text{dead}_A(q)} [\frac{r}{q}, \frac{r+1}{q}) \\
1, & \text{otherwise}
\end{cases}
\end{equation}
\end{definition}

For a Twin prime center ($6N \pm 1$), the forbidden residues are explicitly governed by $6N \equiv \pm 1 \pmod q$. This creates exactly two topological ``holes'' in the phase space. The base survival measure (single-center expectation) is simply the integral over the torus:
\begin{equation}
\mu_q = \int_{0}^{1} \Psi_{\text{Twin},q}(\theta) d\theta = \frac{q-2}{q}
\end{equation}
the complement of the microscopic root $P(q\mid N,\text{twin})=1/(q-2)$ of
Part~XII~\cite{ChenXII}, re-expressed as a geometric transmitting measure.

\section{The Circular Convolution Integral}

We now consider the spatial entanglement of two Twin prime centers separated by distance $j$. The translation from center $N$ to $N+j$ induces a precise phase shift in the $q$-frequency space:
\begin{equation}
\Delta\theta = \frac{j}{q}
\end{equation}

The joint survival measure of both centers is equivalent to the overlapping ``transmitting'' area of their respective indicator functions. This is rigorously expressed as a continuous circular convolution integral on $\mathbb{T}^1$:
\begin{equation}
I_q(j) = (\Psi_q \circledast \Psi_q)\left(\frac{j}{q}\right) = \int_{0}^{1} \Psi_q(\theta) \cdot \Psi_q\left( \left(\theta + \frac{j}{q}\right) \bmod 1 \right) d\theta
\end{equation}

The correlation ratio $R_q(j)$, normalized against the independent probabilities, is therefore:
\begin{equation}
R_q(j) = \frac{I_q(j)}{\mu_q^2} = \frac{q^2}{(q-2)^2} \int_{0}^{1} \Psi_q(\theta) \cdot \Psi_q\left( \theta + \frac{j}{q} \right) d\theta
\end{equation}

\noindent\textbf{Two clarifications.} First, the integral $\int_0^1\Psi_q(\theta)\Psi_q(\theta+s)\,d\theta$ is strictly the \emph{autocorrelation} of $\Psi_q$; it coincides with the circular convolution because the twin's hole set $\{6^{-1},-6^{-1}\}$ is centrally symmetric, so $\Psi_q$ is even. Second, because the lag $j$ is an integer the phase shift $j/q$ is always a multiple of the cell width $1/q$, so this continuous integral is \emph{exactly} the discrete count of co-surviving residues: the phase-space picture is a rigorous geometric reformulation of the singular-series ratio, not an approximation of it.

\begin{figure}[h]
\centering
\includegraphics[width=0.92\textwidth]{fig_circular_convolution.pdf}
\caption{Left: the three interference states of Twin--Twin entanglement on
$\mathbb{T}^1$ for $q=7$. The castellated indicator of center $N$ (blue) and of the
shifted center $N+j$ (red) overlap; grey cells are the $\mathrm{dead}(q)$ holes.
State A ($j\equiv0$) blocks $2/q$, State B ($j\equiv\pm3^{-1}$) blocks $3/q$, State
C (otherwise) blocks the maximal $4/q$. Right: the resulting quantization of
$R_q(j)$ for $q=5,7,11$ into three discrete levels (A,\,B $>1$; C $<1$). At $j=1$
every prime is in State C, so $R(1)=\prod_q R_q(1)=\mathfrak{S}_{\text{quad}}/
\mathfrak{S}_{\text{twin}}^2\approx0.40$.}
\label{fig:convolution}
\end{figure}

The global spatial correlation across the $6N$ lattice is the infinite product of these single-frequency phase-space integrals. Due to the rapid spectral decay of the interference amplitude, this product can be truncated at a spectral cutoff $k^\star(\varepsilon)$ for any desired precision $\varepsilon$:
\begin{equation}
R(j) = \prod_{q > 3}^{k^\star} R_q(j)
\end{equation}

\section{Quantization of Interference States}

Because $\Psi_q(\theta)$ is a step function with two holes of width $1/q$, the continuous integral collapses into three discrete, quantized states, governed entirely by the geometric alignment of the holes. 

The intrinsic phase difference between the two holes of a Twin prime is $\Delta_{\text{holes}} \equiv \pm 2 \cdot 6^{-1} \equiv \pm 3^{-1} \pmod q$.

\begin{theorem}[The Three Interference States]
The circular convolution integral $R_q(j)$ evaluates to exactly three states:
\begin{enumerate}
    \item \textbf{State A (Constructive Interference):} When $j \equiv 0 \pmod q$. The phase shift is $0$. The holes perfectly overlap. Total blocked area is $2/q$.
    $$ R_q(j) = \frac{(q-2)/q}{((q-2)/q)^2} = \mathbf{\frac{q}{q-2}} > 1 $$
    \item \textbf{State B (Partial Interference):} When $j \equiv \pm 3^{-1} \pmod q$. The phase shift aligns one hole of the first center with the other hole of the second. Total blocked area is $3/q$.
    $$ R_q(j) = \frac{(q-3)/q}{((q-2)/q)^2} = \mathbf{\frac{q(q-3)}{(q-2)^2}} > 1 $$
    \item \textbf{State C (Destructive Interference):} For all other $j$. The holes are totally misaligned. Total blocked area is maximal at $4/q$.
    $$ R_q(j) = \frac{(q-4)/q}{((q-2)/q)^2} = \mathbf{\frac{q(q-4)}{(q-2)^2}} < 1 $$
\end{enumerate}
\end{theorem}

\vspace{0.5em}
\noindent \textbf{Remark on Spectral Decay:} 
Observe that across all three states, as the spatial frequency $q \to \infty$, the correlation ratio asymptotically approaches unity: $\lim_{q \to \infty} R_q(j) = 1$. This rigorously proves that high-frequency topological holes become vanishingly small, and the spatial interference is overwhelmingly dominated by the lowest frequencies (small primes). This analytic property justifies the existence of the spectral cutoff $k^\star(\varepsilon)$ defined in Equation (6) and guarantees the absolute convergence of the global lattice correlation product.

Quantitatively, the per-prime deviation from unity is $|\ln R_q|\approx 2/q$ for State A, $\approx 1/q$ for State B, and $\approx 4/(q-2)^2=O(1/q^2)$ for State C. For any fixed lag $j$, every prime $q>3j+1$ is in State C (only the finitely many primes dividing $j$ or $3j\pm1$ can be A or B), so the truncation tail is governed by the convergent sum $\sum_{q>k^\star}4/(q-2)^2$. Evaluating the product numerically for the canonical case $j=1$ (all State C) gives the cutoff $k^\star(\varepsilon)$ at which the truncated product is within relative error $\varepsilon$ of the full product $R(1)=0.3969$:

\begin{table}[h]\centering
\begin{tabular}{ccc}
\hline
$\varepsilon$ & $k^\star$ & primes retained\\
\hline
$10^{-1}$ & $13$ & $4$\\
$10^{-2}$ & $73$ & $19$\\
$10^{-3}$ & $557$ & $100$\\
$10^{-4}$ & $4229$ & $577$\\
\hline
\end{tabular}
\caption{Spectral cutoff $k^\star(\varepsilon)$ for the truncated lattice
correlation product at $j=1$. To $1\%$ precision only the primes up to $73$ are
needed; the $O(1/q^2)$ State-C tail makes the product converge rapidly.}
\label{tab:kstar}
\end{table}

\section{Limit Case Proof: The Quadruplet Collapse}

To demonstrate the analytic power of this geometric formulation, we evaluate the limit case $j=1$, which physically corresponds to adjacent Twin centers merging into a Prime Quadruplet $(6N-1, 6N+1, 6N+5, 6N+7)$.

Let us test the spatial frequency $q=5$. We compute the phase difference:
$3^{-1} \pmod 5 \equiv 2$.
Since $j=1$, we have $1 \not\equiv 0 \pmod 5$ and $1 \not\equiv \pm 2 \pmod 5$. 

Thus, $j=1$ deterministically triggers \textbf{State C (Destructive Interference)}. Substituting into the State C equation yields:
\begin{equation}
R_5(1) = \frac{5(5-4)}{(5-2)^2} = \mathbf{\frac{5}{9}} \approx 0.555
\end{equation}

This is the largest single-prime suppression, and the product over all primes gives
$R(1)=\prod_{q>3}R_q(1)=\mathfrak{S}_{\text{quad}}/\mathfrak{S}_{\text{twin}}^2\approx0.40$:
the quadruplet is about $0.40\times$ as common as two independent twins. We are
careful about what this proves. It is the suppression of the quadruplet
\emph{prefactor}---its $X$-independent singular series relative to the product of
its constituents---and it is a deterministic consequence of four non-overlapping
phase-space holes maximizing destructive interference on the integration manifold.

It is \emph{not}, however, the $-4$ proper-depth slope of Volume~I. The geodynamic
law $\ln\rho_4=\ln\mathfrak{S}_4-4\,\ell$ separates cleanly into an intercept
$\ln\mathfrak{S}_4$ and a slope $-4$; the destructive interference computed here
fixes the intercept (through $\mathfrak{S}_4$), while the slope $-4$ is the bare
member count, the four factors of $1/\ln X$, independent of any inter-twin
correlation. The two mechanisms are complementary and live on the orthogonal axes
of \S1: interference explains \emph{why the constant is small}, dimension explains
\emph{why the density decays}. Conflating them would re-merge the very axes the
framework is built to separate.

\section{Conclusion}
By translating modular arithmetic into the circular convolution of castellated indicator functions, we have established a pure continuous-integral foundation for spatial entanglement on the 6N lattice---the analytic form of the quasi-periodic congruence diffraction pattern identified empirically in the Volume~I Review~\cite{Review}. Future work will extend this 1D phase-space integration into multidimensional tensor products to compute the absolute spectral cutoff of the congruence diffraction pattern.

\paragraph{Data and code availability.}
The figure-generating script and an independent verifier (the three-state formula
against direct union counts, and the $k^\star(\varepsilon)$ error table) are openly
available at \url{https://github.com/Ruqing1963/6N-phase-space}~\cite{repo}.

\begin{thebibliography}{9}
\bibitem{HL} G.~H.~Hardy, J.~E.~Littlewood, \emph{Some problems of `Partitio
Numerorum'; III}, Acta Math.\ \textbf{44} (1923), 1--70.
\bibitem{ChenI} R.~Chen, \emph{Prime distribution on the $6N$ skeleton} (Part~I),
Zenodo, doi:10.5281/zenodo.20470367 (2026).
\bibitem{ChenXII} R.~Chen, \emph{The microscopic root $1/(q-2)$ and the mean
shift} (Part~XII), Zenodo, doi:10.5281/zenodo.20528446 (2026).
\bibitem{ChenXVII} R.~Chen, \emph{Arithmetic geodynamics on the $6N$ skeleton}
(Part~XVII), Zenodo, doi:10.5281/zenodo.20541575 (2026).
\bibitem{ChenXIX} R.~Chen, \emph{The tuple-size law at $m=5$: resolving the sparse
regime} (Part~XIX), Zenodo, doi:10.5281/zenodo.20551393 (2026).
\bibitem{Review} R.~Chen, \emph{Arithmetic geodynamics on the $6N$ skeleton: a
review of Parts I--XIX}, \url{https://github.com/Ruqing1963/6N-review} (2026).
\bibitem{repo} R.~Chen, \emph{The circular convolution of prime constellations in
congruence phase space: code} (Part~XX),
\url{https://github.com/Ruqing1963/6N-phase-space} (2026).
\end{thebibliography}

\end{document}
